\documentclass[UTF8,a4paper,11pt,fontset=windows]{ctexart}
\usepackage[left=19mm,right=19mm,top=17mm,bottom=18mm]{geometry}
\usepackage{amsmath,amssymb,array,longtable,booktabs,needspace,fancyhdr}
\usepackage[hidelinks]{hyperref}
\setlength{\parindent}{0pt}
\setlength{\parskip}{4pt}
\linespread{1.13}
\pagestyle{fancy}\fancyhf{}\renewcommand{\headrulewidth}{0pt}
\fancyfoot[C]{\small 第\thepage 页}
\setlength{\headheight}{14pt}
\newcommand{\R}{\mathbb R}
\newcommand{\src}[1]{\par{\scriptsize 教材来源：#1。}\par}
\newcommand{\q}[2]{\par\Needspace{4\baselineskip}\textbf{#1.}\ #2}
\newcommand{\subq}[2]{\par\Needspace{2\baselineskip}（#1）#2}
\newcommand{\opts}[4]{\par\begin{tabular*}{\linewidth}{@{}p{.475\linewidth}@{\extracolsep{\fill}}p{.475\linewidth}@{}}A. #1 & B. #2 \\ C. #3 & D. #4\end{tabular*}\par}
\newcommand{\optsfour}[4]{\par\begin{tabular*}{\linewidth}{@{}p{.235\linewidth}@{\extracolsep{\fill}}p{.235\linewidth}p{.235\linewidth}p{.235\linewidth}@{}}A. #1 & B. #2 & C. #3 & D. #4\end{tabular*}\par}
\newcommand{\parttitle}[1]{\par\Needspace{4\baselineskip}{\bfseries #1}\par}
\newcommand{\work}[1]{\par\vspace*{#1}}
\newcommand{\pts}[1]{\unskip\nobreak\hfill\mbox{（#1分）}\par}
\newcommand{\ans}[2]{\par\Needspace{3\baselineskip}\textbf{第#1题}\quad #2\par}
\newcommand{\titleblock}[2]{\begin{center}{\LARGE\bfseries #1}\par\vspace{3mm}{\normalsize #2}\end{center}}
\renewcommand{\arraystretch}{1.35}

\begin{document}
\titleblock{高一数学期中教材改编练习卷教师版}{答案解析\qquad 评分建议\qquad 教材出处\qquad 复习诊断}
\parttitle{范围与结构}
本卷供新高一学生期中限时训练使用。依据学校命题思路表，范围为湘教版必修第一册第一章至第四章第2节指数函数，不涉及对数函数。共19题，120分钟，150分。所有题目均由所提供教材的例题、练习或习题改编；出处和改动见第6页。

\begin{tabular*}{\linewidth}{@{\extracolsep{\fill}}lrrrrr@{}}
\toprule 模块 & 单选 & 多选 & 填空 & 解答 & 合计\\\midrule
集合与逻辑 & 15 & 0 & 0 & 15 & 30\\
一元二次函数、方程和不等式 & 5 & 6 & 5 & 15 & 31\\
函数的概念与性质 & 20 & 12 & 10 & 17 & 59\\
幂函数、指数函数 & 0 & 0 & 0 & 30 & 30\\\midrule
合计 & 40 & 18 & 15 & 77 & 150\\\bottomrule
\end{tabular*}

\textbf{与命题表的对应。}集合：第1、15题，共20分；逻辑：第2、3题，共10分。基本不等式：第4题5分；不等式性质：第9题6分；恒成立：第12题5分；二次不等式：第16题15分。函数概念与表示：第5、6、14题15分；函数性质：第7、8、10、11、13、17题44分。第18题计入幂函数及实数指数幂17分，第19题为指数函数13分。

\textbf{分值次序。}解答题依内容安排为15、15、17、17、13分，与命题表各模块的解答题分值保持一致。练习采用8道单选、3道多选、3道填空、5道解答的结构；难度按新高一教材复习设计，不把高考试题难度套入未学内容。

\parttitle{答案速查}
\begin{tabular*}{\linewidth}{@{\extracolsep{\fill}}lrrrrrrrr@{}}
\toprule 题号 & 1 & 2 & 3 & 4 & 5 & 6 & 7 & 8\\\midrule
答案 & B & A & C & B & C & B & D & A\\\bottomrule
\end{tabular*}
\begin{tabular*}{\linewidth}{@{\extracolsep{\fill}}lccc@{}}
\toprule 题号 & 9 & 10 & 11\\\midrule
答案 & AC & AB & ABD\\\bottomrule
\end{tabular*}

第12题：$[1,+\infty)$。\quad 第13题：$f(-3)<f(2)<f(-1)$。

第14题：$[-2,2]$。

\textbf{评分约定。}单选、填空答对各得5分。多选全部选对6分，选对非空真子集且无错误选项3分，空选或含错误选项0分。此为本练习统一约定，如学校另有细则可调整。解答题下列分值为建议，等价正确解法照常给分；不因重复使用同一结论重复扣分。

\textbf{难度说明。}第7页的“基础、常规、中档”是备课判断，尚未经学生试测。不能据此预估班级均分、达标率或考试成绩。

\newpage
\parttitle{第1至14题解析}
\ans{1}{B。两集合的公共元素为2、4，故$A\cap B=\{2,4\}$。}
\ans{2}{A。$x=2\Rightarrow x^2=4$；反向不成立，如$x=-2$。}
\ans{3}{C。将“任意”改为“存在”，并将“$\ge0$”改为“$<0$”。命题的否定为“$\exists x\in\R$，$x^2-4x+4<0$”。}
\ans{4}{B。因$x-2>0$，有
\[
x+\frac4{x-2}=2+(x-2)+\frac4{x-2}\ge2+2\sqrt4=6.
\]
等号成立当且仅当$x-2=\frac4{x-2}$，结合$x>2$得$x=4$，符合题设。}
\ans{5}{C。$\sqrt{x^2}=|x|$，两者定义域都是$\R$。A中第二个函数在$x=-1$处无定义；B的对应关系不同；D在负数处的函数值不同。}
\ans{6}{B。由$x+2\ge0$且$(x-3)(x+2)\ne0$，得$x>-2$且$x\ne3$。定义域为$(-2,3)\cup(3,+\infty)$。}
\ans{7}{D。$f(2)=4+\frac12=\frac92$，奇函数满足$f(-2)=-f(2)=-\frac92$。}
\ans{8}{A。对称轴为$x=1-a$。在$(-\infty,2]$上单调递减当且仅当$1-a\ge2$，即$a\le-1$。对称轴可以恰好位于区间右端点。}
\ans{9}{AC。$a<b<0$，故$-a>-b$。又$ab>0$，
$\frac1a-\frac1b=\frac{b-a}{ab}>0$，B错误。
$\sqrt{a^2}=|a|=-a>-b=|b|$，C正确。由$a<b<0$两边除以$a$，得$1>\frac ba>0$，D错误。}
\ans{10}{AB。$f(x)=(x-2)^2-3$，在$[0,2]$上递减，在$[2,3]$上递增。
最小值为$f(2)=-3$；$f(0)=1$，$f(3)=-2$，最大值为1，值域是$[-3,1]$。}
\ans{11}{ABD。令$x=y=0$，得$f(0)=2f(0)$，所以$f(0)=0$。
令$y=-x$，得$f(-x)=-f(x)$，故为奇函数。
$f(3)=f(2)+f(1)=3f(1)=6$。若同时为偶函数，则$f(1)=f(-1)=-f(1)$，与$f(1)=2$矛盾，C错误。无须推断$f$对任意实数均为线性函数。}
\ans{12}{$x^2-2x+m=(x-1)^2+m-1$，其在$\R$上的最小值为$m-1$。恒成立等价于$m-1\ge0$，故$m\in[1,+\infty)$。}
\ans{13}{偶性给出$f(-3)=f(3)$，$f(-1)=f(1)$；因$3>2>1\ge0$且在非负半轴递减，故$f(-3)<f(2)<f(-1)$。}
\ans{14}{去绝对值后，
\[
f(x)=\begin{cases}2,&x<-1,\\-2x,&-1\le x\le1,\\-2,&x>1.\end{cases}
\]
中间一段可取遍$[-2,2]$，两端值也能取到，故值域为$[-2,2]$。}

\newpage
\parttitle{第15与16题解析及评分}
\ans{15}{（15分）}
\subq{1}{当$a=-1$时，$A=(-\infty,2]$，$B=(-1,+\infty)$。\pts{2}}
于是
\[
A\cap B=(-1,2],\qquad A\cup B=\R,\qquad \complement_U A=(2,+\infty).
\]
三个结果各2分。\pts{6}
\subq{2}{若$a\le2$，任取$x\in\R$：当$x\le2$时，$x\in A$；当$x>2$时，$x>a$，故$x\in B$。因此$A\cup B=\R$。\pts{3}}
若$a>2$，则$x=a$既不属于$A$，也不属于$B$，所以$A\cup B\ne\R$。\pts{3}
综上，$a$的取值范围为$(-\infty,2]$。\pts{1}

\textbf{边界核验。}当$a=2$时，$A$包含2，而$B$包含一切大于2的实数，二者的并集仍为$\R$，不能漏掉等号。第（1）问共8分，第（2）问共7分。

\ans{16}{（15分）}
\subq{1}{因$2x^2-x-3=(2x-3)(x+1)$，\pts{2}}
对应方程的两个根为$-1$、$\frac32$。\pts{2}
二次项系数为正，故不等式的解集为
\[
\boxed{\left[-1,\frac32\right]}.
\]
写出两根之间并正确保留端点。\pts{3}
\subq{2}{原不等式可化为$(x-m)(x-2)<0$。\pts{2}}
当$m<2$时，解集为$(m,2)$；\pts{2}

当$m=2$时，原不等式为$(x-2)^2<0$，解集为$\varnothing$；\pts{2}

当$m>2$时，解集为$(2,m)$。\pts{2}

\textbf{评分提醒。}第（1）问共7分，第（2）问共8分。第（2）问中的$m=2$是独立的一类；不能在未比较两根大小时统一写成$(m,2)$。写成等价的完整分类形式可给满分。

\newpage
\parttitle{第17与18题解析及评分}
\ans{17}{（17分）}
\subq{1}{$f$的定义域为$[-3,3]$，关于原点对称。\pts{2}}
任取$x\in[-3,3]$，有$f(-x)=g(-x)+g(x)=f(x)$，因此$f$为偶函数。\pts{3}
\subq{2}{当$x\in[-3,0)$时，$-x\in(0,3]$。\pts{2}}
由偶性，$f(x)=f(-x)=(-x)^2-2(-x)=x^2+2x$。\pts{2}
\subq{3}{在$[0,3]$上，$f(x)=(x-1)^2-1$，所以最小值为$-1$，在$x=1$时取得。\pts{2}}
端点值为$f(0)=0$，$f(3)=3$，最大值为3，在$x=3$时取得。\pts{2}
由偶性可知整个定义域内，最小值为$-1$，对应$x=\pm1$；最大值为3，对应$x=\pm3$。\pts{4}

\textbf{核验。}全式为$f(x)=x^2-2|x|=(|x|-1)^2-1$，$|x|\in[0,3]$。这也可独立确认全域最值及其全部取值点。无需确定$g$本身；例如$g=f/2$能满足题设，条件彼此相容。

\ans{18}{（17分）}
\subq{1}{$27^{\frac23}=9$，$16^{-\frac34}=\frac18$，\pts{3}}
所以$27^{\frac23}+16^{-\frac34}=\frac{73}{8}$。\pts{2}
\subq{2}{由$2^\alpha=\frac14=2^{-2}$，得$\alpha=-2$，故$f(x)=x^{-2}=\frac1{x^2}$。\pts{3}}
定义域为$(-\infty,0)\cup(0,+\infty)$。\pts{1}
\subq{3}{定义域关于原点对称，且$f(-x)=\frac1{(-x)^2}=f(x)$，故为偶函数。\pts{2}}
在$(-\infty,0)$上单调递增，在$(0,+\infty)$上单调递减。\pts{2}
下面用定义证明负半轴的单调性。
任取$x_1<x_2<0$，有
\[
f(x_2)-f(x_1)=\frac1{x_2^2}-\frac1{x_1^2}
=\frac{(x_1-x_2)(x_1+x_2)}{x_1^2x_2^2}>0,
\]
因为$x_1-x_2<0$，$x_1+x_2<0$，分母大于0，故$f(x_1)<f(x_2)$。因此$f$在$(-\infty,0)$上单调递增。\pts{4}

\textbf{评分提醒。}第（1）问5分，第（2）问4分，第（3）问8分。不能把两个分离的单调区间合写成“在定义域上单调”；只写结论、没有定义证明，不得相应证明分。

\newpage
\parttitle{第19题解析及评分}
\ans{19}{（13分）}
\subq{1}{由$f(1)=a^{|2-1|}=a=\frac14$，得$a=\frac14$。\pts{3}}
\subq{2}{因为$0<\frac14<1$，函数$y=(\frac14)^t$关于$t$单调递减。}
当$x\le\frac12$时，$|2x-1|=1-2x$随$x$增大而减小，因此$f$在$(-\infty,\frac12]$上单调递增。\pts{2}

当$x\ge\frac12$时，$|2x-1|=2x-1$随$x$增大而增大，因此$f$在$[\frac12,+\infty)$上单调递减。\pts{2}

又$|2x-1|$可取遍$[0,+\infty)$，故$f$的值域为$(0,1]$，最大值1在$x=\frac12$处取得，0不能取得。\pts{3}
\subq{3}{因$\frac1{16}=(\frac14)^2$且底数在$(0,1)$内，}
\[
\left(\frac14\right)^{|2x-1|}\ge\left(\frac14\right)^2
\iff |2x-1|\le2.
\]
同底比较时正确改变指数间不等号方向。\pts{1}
解得$-2\le2x-1\le2$，即$-\frac12\le x\le\frac32$。\pts{1}
不等式的解集为$\boxed{[-\frac12,\frac32]}$。\pts{1}

\textbf{核验。}代入两端点均得$f(x)=\frac1{16}$，端点应保留；代入$x=\frac12$得$f(x)=1$，满足不等式。第（1）问3分，第（2）问7分，第（3）问3分。

\parttitle{练后反馈与回到教材的方法}
建议先完整限时完成学生卷，再依据实际作答记录定位问题。参考时间为：单选25分钟、多选15分钟、填空10分钟、解答60分钟、检查10分钟，共120分钟。这是练习安排建议，不是学校规定。

第6、15题若出现端点错误，要求学生在数轴上分别标出根式条件、分母零点与集合端点，补做教材p.13第9、16题及p.77第4题。第16题漏掉$m=2$，则先列出$m<2$、$m=2$、$m>2$的根的位置，再分别画符号表。

第13、17题若把负半轴直接套入正半轴的增减结论，先用偶性把比较转到非负半轴，再说明比较方向。第18题证明中若符号判断错误，应明确写出两个负因子和正分母；第19题若出现方向或值域端点错误，则分别核对底数小于1、指数非负、指数最小值能否取得。

学生纠错时可记录“我的错误等式或判断”“教材依据”“修正后的关键一步”，再做一次只改数字的同类题。上述卡点均为备课预判，须用本次练习的真实作答验证。

\newpage
\parttitle{逐题教材出处与改编记录}
下表“页”为教材印刷页码；所提供PDF的文件页序为印刷页码加7，例如p.13对应PDF第20页。题号以原页图像核对，未用乱码文本推测题号。来源均为同一本所提供教材。
{\small
\begin{longtable}{@{}p{.05\linewidth}p{.32\linewidth}p{.56\linewidth}@{}}
\toprule 题 & 原题定位 & 本卷的改动\\\midrule\endhead
1 & p.13 习题1.1第8题 & 改两个有限集的元素，仅保留交集，改成单选。\\
2 & p.18 例5（1） & 将$x=1$与$x^2=1$改成$x=2$与$x^2=4$，改成条件判断单选。\\
3 & p.23 页首练习（1） & 将$x^2-2x+1$改为$x^2-4x+4$，保留全称命题否定，改成单选。\\
4 & p.43 习题2.1第7（1）题 & 将$x>1$、$x+1/(x-1)$改为$x>2$、$x+4/(x-2)$，保留配凑最值结构。\\
5 & p.77 习题3.1第5题 & 保留定义域与对应关系两项判据，重组选项，正确项改为$\sqrt{x^2}$与$|x|$。\\
6 & p.77 习题3.1第4（1）题 & 根式改为$\sqrt{x+2}$，分母改为$x^2-x-6$，改成定义域单选。\\
7 & p.86 习题3.2第6题 & 保留正半轴解析式，将求$f(-1)$改为求$f(-2)$，明确完整定义域，改成单选。\\
8 & p.82 例3 & 仅将递减区间右端点由4改为2，改成单选。\\
9 & p.43 习题2.1第3题 & 保留$a<b<0$及四类比较，将填符号改为多选判断。\\
10 & p.82 练习第2（3）题；p.90 复习题三第7题 & 改二次函数系数及闭区间，合并单调性、最大最小值和值域为多选。\\
11 & p.91 复习题三第12题 & 保留加法关系，新增$f(1)=2$，将数值结论改为$f(3)=6$；未假定连续性。\\
12 & p.55 例5 & 将原含参二次式简化为$x^2-2x+m$，严格大于改为大于等于，改成填空。\\
13 & p.86 习题3.2第11题 & 保留偶性与正半轴递减条件，比较对象改为$f(-3)$、$f(2)$、$f(-1)$。\\
14 & p.76 页首练习第1（2）题 & 将$|x-2|-|x+2|$改为$|x-1|-|x+1|$，只问值域。\\
15 & p.13 习题1.1第9、16题 & 合并区间运算和并集覆盖问题，取$A=(-\infty,2]$、$B=(a,+\infty)$，分为数值与参数两问。\\
16 & p.59 习题2.3第4、8题 & 第（1）问改为$2x^2-x-3\le0$；第（2）问将相邻根改为$m$、2，引入根的顺序与重根分类。\\
17 & p.84 例5；p.82 练习第2（3）题 & 保留$g(x)+g(-x)$构造，区间改为$[-3,3]$，非负段改为$x^2-2x$，增设全域最值与全部取值点。\\
18 & p.104 习题4.1第1、4、8题 & 合并实数幂运算与幂函数过点，数值改为$27^{2/3}+16^{-3/4}$和$(2,1/4)$；作图改为定义证明。\\
19 & p.114 习题4.2第11题 & 将$|3x-1|$改为$|2x-1|$，$f(1)$改为$1/4$，保留底数、单调性、值域，增设同底比较。\\
\bottomrule
\end{longtable}}

\newpage
\parttitle{逐题教学目标与诊断建议}
难度为备课预判。“观察点”用于批改和反馈，不代表已有学生表现。选择、填空题按第1页规则计分；解答题按第3至5页给分。
{\small\setlength{\tabcolsep}{3pt}
\begin{longtable}{@{}p{.04\linewidth}>{\raggedright\arraybackslash}p{.17\linewidth}>{\raggedright\arraybackslash}p{.22\linewidth}>{\raggedright\arraybackslash}p{.28\linewidth}>{\raggedright\arraybackslash}p{.17\linewidth}@{}}
\toprule 题 & 目标与难度 & 预期方法 & 可能错误 & 评价观察点\\\midrule\endhead
1 & 交集；基础 & 寻找公共元素 & 交集当并集 & 是否只取公共元素\\
2 & 充分必要；基础 & 正推并举反例 & 反向推理漏$-2$ & 能否说明反向不成立\\
3 & 命题否定；基础 & 换量词与否定结论 & 只换量词；把$\ge$否定为$\le$ & 量词和关系符号同时正确\\
4 & 基本不等式；常规 & 令$t=x-2>0$配定积 & 直接把$x$与分式配积；忽略取等 & 下界和可取性均核实\\
5 & 同一函数；基础 & 同时比定义域和对应关系 & 只看化简后的式子 & 能否指出其余组的差别\\
6 & 定义域；常规 & 联立根式与分母条件 & 忘排除$x=-2$或3 & 两类限制无遗漏\\
7 & 奇性求值；基础 & 先算正值再变号 & 把$-2$直接代入正段 & 输入所在区间判断正确\\
8 & 单调性含参；常规 & 对称轴与区间右端比较 & 写成等于；漏边界 & $1-a\ge2$方向与等号\\
9 & 不等式性质；常规 & 作差；除负数反向；绝对值 & 倒数顺序错；$\sqrt{a^2}=a$ & 四项逐一有依据\\
10 & 闭区间最值；基础 & 配方并比较端点 & 最大值只看右端点 & 顶点及两个端点齐全\\
11 & 奇偶性推理；中档 & 代0、代相反数、逐次相加 & 无条件认定所有实数上$f(x)=2x$ & 有限次代入足以支撑结论\\
12 & 恒成立；常规 & 配方求全域最小值 & 把任意误为存在；漏$m=1$ & 最小值非负的充要性\\
13 & 偶性与单调性；常规 & 负输入转正后比较 & 负半轴直接沿用递减 & 比较对象在同一单调区间\\
14 & 值域与表示；常规 & 按$-1$、1分段 & 仅报上下界；漏端点 & 中间段能取遍全区间\\
15 & 区间运算；常规 & 数轴及覆盖检验 & 漏$a=2$；补集端点错 & 三个集合与参数边界\\
16 & 二次不等式；中档 & 因式分解后分类根序 & 将$(m,2)$当统一答案 & 三种参数情形完整\\
17 & 偶性与全域最值；中档 & 定义证明、负段还原、配方 & 负段错号；最值点只报正数 & 定义域、证明、全体最值点\\
18 & 实数幂与幂函数；中档 & 运算法则、过点求幂、定义作差 & 负指数错；漏0的限制；负段符号错 & 四个因子符号及严格不等式\\
19 & 指数函数；中档 & 先定底数再按绝对值分段 & 底数小于1仍同向；把0收入值域 & 增减、值域、闭端点齐全\\
\bottomrule
\end{longtable}}

\textbf{复核记录。}逐题以原页图像确认来源，重算所有答案并核对选项；额外检查第4题取等、第6题分母零点、第8题对称轴边界、第15题$a=2$、第16题$m=2$、第17题条件相容及最值点、第18题负半轴定义证明、第19题指数比较方向与值域端点。总分与各模块分值已核对。正式使用前可按班级进度调整训练节奏。
\end{document}
